\section{Signed cancellation along the source trajectories}
\label{can:section}

The direct force estimates lose an angular factor when the source velocity
is much closer to the cone direction than the receiver velocity is. In
this region we integrate along the source trajectory before taking absolute
values. The transport part of the force must be included: it cancels a
term produced by differentiating the cone geometry.

Keep a supported receiver $X(t)$, with velocity $a=X'$, energy $q_X$ and
momentum derivative $K=V_X'$. A source label has position $Y(s)$, velocity
$u=Y'$ and energy $q$. Along its retarded branch, set
\[
 t=s+r,\quad rn=X(t)-Y(s),\quad
 d=1-n\cdot u,\quad D=1-n\cdot a,\quad e=1-a\cdot u,
 \qquad |n|=1.
\]
A prime denotes differentiation in $s$ along this branch, and $a_t$ denotes
the receiver derivative in its own time. By
Lemma~\ref{ret:receiver-time}, the retarded measure identities are
\begin{equation}\label{can:measure}
 \frac{dt}{ds}=\frac dD,
 \qquad ds\,f_0(z)\,dz=D\,dt\,f(s,y,v)\,dy\,dv.
\end{equation}
In particular $d,D>0$. Differentiating the retarded relation gives
\begin{equation}\label{can:geometry}
 r'=\frac dD-1,\qquad
 rn'=N_0:=\frac dD(a-n)+n-u,\qquad n\cdot N_0=0.
\end{equation}
Define the matrices and vectors
\begin{equation}\label{can:definitions}
 Q=D\Id+n\otimes a,\qquad H=Q/D,\qquad
 g=\frac{u-n}{d},\qquad k_0=u-\frac eD n.
\end{equation}
Thus $Hg=k_0/d$.

\subsection{The exact identity}

After the change to $ds\,f_0(z)dz$, the two interior force terms in
\eqref{ret:force-kernels} together have kernel
\begin{equation}\label{can:input-kernel}
 \mathcal K=-\frac Hr\,\partial_s^{(u)}g
             -\frac{Hg}{r^2q^2d}.
\end{equation}
Here $\partial_s^{(u)}$ differentiates through $u$ alone. The second term
is the transport kernel. As in the force representation, we suppress the
common constant $1/(4\pi)$ throughout this section.

\begin{lemma}[Signed force identity]\label{can:identity}
At every non-tip point of a retarded branch,
\begin{equation}\label{can:label-identity}
 \mathcal K
 =-\left(\frac{k_0}{rd}\right)'
  +\frac{e}{r^2D^2}\left(\frac{n}{q_X^2D}-a\right)
  +\frac{n(a_t\cdot k_0)}{rD^2}.
\end{equation}
Consequently the two nondifferentiated terms, in $dt\,dy\,dv$ measure,
are
\begin{equation}\label{can:bulk}
 \mathcal R=\frac{fe}{r^2D}\left(\frac{n}{q_X^2D}-a\right),
 \qquad
 \mathcal A=\frac{fn(a_t\cdot k_0)}{rD}.
\end{equation}
\end{lemma}

\begin{proof}
The product rule first rewrites the source-acceleration term as
\[
 -\frac Hr\partial_s^{(u)}g
 =-(Hg/r)'+(H/r)'g+(H/r)\partial_s^{(n)}g.
\]
For a variation $A$ of $a$, with $u,n$ fixed,
\[
 \partial_a k_0[A]=\frac nD(A\cdot k_0).
\]
Since $a'=a_t d/D$, its contribution is the last term of
\eqref{can:label-identity}.
The terms from $r,n$, including the transport term, are
\[
 \frac1r\left(\frac{k_0}{d}\right)'_n
       -\frac{k_0r'}{r^2d}-\frac{k_0}{r^2q^2d^2}.
\]
Multiplying by $r^2d^2$ gives
\[
 rd(k_0)'_n+k_0(N_0\cdot u-dr'-q^{-2}).
\]
Directly from \eqref{can:geometry} and $|u|^2=1-q^{-2}$,
\[
 N_0\cdot u-dr'-q^{-2}=-\frac{ed}{D},
 \qquad
 rd(k_0)'_n=-\frac{ed}{D}
            \left(N_0+\frac{n(N_0\cdot a)}D\right).
\]
Finally, $|a|^2=1-q_X^{-2}$ gives
\[
 N_0+\frac{n(N_0\cdot a)}D+k_0
   =\frac dD\left(a-\frac{n}{q_X^2D}\right).
\]
These identities prove the middle term in
\eqref{can:label-identity}, with its stated sign. Multiplication by the
factor $fD$ in \eqref{can:measure} gives \eqref{can:bulk}.
\end{proof}

The two residual densities $\mathcal R$ and $\mathcal A$ contain neither
the source acceleration $b$ nor inverse powers of the source deficit $d$.
The geometric residual $\mathcal R$ will be estimated directly. The
receiver term $\mathcal A$ contains $a_t$, so its coefficient must be
small enough to use an absolute bound on the receiver force. The
primitive $k_0/(rd)$ retains the source singularity: after integration
by parts it creates boundary and cutoff-derivative terms. We now keep
track of those terms.

\subsection{Smooth cutoffs and the small angular sector}

Use the dyadic representatives $p,\theta,\phi,h$ for, respectively,
\[
 q,\qquad \sqrt{d/2},\qquad \sqrt{D/2},\qquad r/\max(1,T_0).
\]
The fixed factor in the convention for $p$ is retained. Thus
$q\asymp p$, $q_X\asymp w$, $d\asymp\theta^2$,
$D\asymp\phi^2$ and $r\asymp h$, with constants depending on $T_0$
where necessary. The small-sector representatives satisfy
$\theta<\kappa\phi$. All estimates below remain valid on the wider
overlap range $\theta\leq16\kappa\phi$, after choosing the fixed
$\kappa>0$ sufficiently small.

Use the smooth partitions chosen in Section~\ref{occ:section}, whose
factors have uniformly bounded derivatives in the logarithm of each
variable. For an index $j=(p,\theta,\phi,h)$, let $\beta_j$ be the
product of its four factors. These products satisfy
\begin{equation}\label{can:partition}
 \beta_j\geq0,\qquad \sum_j\beta_j=1.
\end{equation}
The sum is locally finite for $r>0$, with uniformly bounded overlap.
Derivative bounds are taken on the closed support indicators; no estimate
of a derivative by the bump itself is assumed.

\begin{lemma}[Derivatives of a selected sum]\label{can:cutoff-cancellation}
Let $\mathcal S$ be any fixed set of representative indices, all in the
small sector, and put $\chi=\sum_{j\in\mathcal S}\beta_j$. Selection
may depend on the interval and fixed parameters, but not on the variables
being differentiated. For differentiation through source variables
$(q,u)$ alone, or through geometric variables $(r,n)$ alone,
\begin{equation}\label{can:partition-cancel}
 \partial\chi=-\sum_{j\notin\mathcal S}\partial\beta_j.
\end{equation}
This derivative vanishes unless selected and unselected supports meet.
On their intersection, representatives in each coordinate differ by at
most a factor of four. In particular all relevant supports remain in
the wider small-sector range above.
\end{lemma}

\begin{proof}
Differentiate the locally finite identity \eqref{can:partition} before
restricting to a retarded trajectory. This proves
\eqref{can:partition-cancel} separately for each indicated partial
derivative. A smooth nonnegative function and its first derivatives vanish
at every point where the function vanishes. Thus if either family has no
positive summand, the corresponding derivative sum is zero. Meeting
supports constrain each representative to within a factor of two of the
same variable, proving the factor-four comparison. The ratio $\theta/\phi$
therefore changes by at most a factor of sixteen.
\end{proof}

The cancellation applies to the cutoff derivative only after its common
physical multiplier has been factored out. It does not compare derivatives
of two different kernels. An additional $\sigma$-partition may subsequently
be inserted into a nonnegative remaining-sector bound at the structural
edge; it is not differentiated in \eqref{can:partition-cancel}.

\subsection{Bounds for each derivative}

The Lorentz acceleration satisfies
\[
 a_t=(\Id-a\otimes a)K/q_X.
\]
The source estimates needed here, with
$G=|E+n\times B|+|n\cdot B|$ at the source point, are
\begin{equation}\label{can:source-input}
 |u'|\lesssim p^{-1}(G+\theta^2|B|),\quad
 |n\cdot u'|\lesssim p^{-1}\theta(G+\theta^2|B|),\quad
 |q'|\lesssim G+\theta|B|.
\end{equation}
These are \eqref{ret:acceleration-bounds} with $q\asymp p$.
For nonzero $V_X$, let $P_X$ denote orthogonal projection perpendicular
to $V_X$. In the small sector,
\begin{equation}\label{can:vector-bounds}
 |e-D|\lesssim\phi\theta,\quad
 |k_0|\lesssim\theta/\phi,\quad
 |P_Xk_0|\lesssim\theta,\quad
 |P_Xn|\lesssim\phi.
\end{equation}
Indeed $e-D=a\cdot(n-u)$; splitting $a=n+(a-n)$ proves the first bound,
and then $k_0=(u-n)+n(1-e/D)$ proves the next two.
Also $P_Xa=0$. Since $q_X^{-2}\lesssim D$,
\begin{equation}\label{can:bulk-bounds}
 |\mathcal R|\lesssim f/r^2,\quad
 |P_X\mathcal R|\lesssim f\phi/r^2,\quad
 |\mathcal A|\lesssim\frac{f\theta|K|}{wr\phi^2},\quad
 |P_X\mathcal A|\lesssim\frac{f\theta|K|}{wr\phi}.
\end{equation}
For the last two bounds, $|a_t|\lesssim|K|/w$ and
$|n\cdot a_t|\lesssim\phi|K|/w$. Insert the decomposition of $k_0$
in \eqref{can:vector-bounds} to obtain
$|a_t\cdot k_0|\lesssim\theta|K|/w$.

After integration by parts, a cutoff derivative in source time has common
multiplier $fDk_0/(rd)$ in $dt\,dy\,dv$ measure. Its size is at most
$Cf\phi/(r\theta)$, with an additional factor $\phi$ after projection.
The separate logarithmic rates and resulting bounds are as follows.

\paragraph{Source derivatives.}
From \eqref{can:source-input},
\[
 \frac{|q'|}{q}+\frac{|d'_u|}{d}
 \lesssim\frac1p\left(\frac G\theta+\theta|B|\right).
\]
Thus their error kernels are bounded by
\begin{equation}\label{can:source-cutoff-bound}
 C\left(\frac{f\phi}{rp\theta^2}G
            +\frac{f\phi}{rp}|B|\right).
\end{equation}
They coincide with the direct source-force majorants. By
Lemma~\ref{can:cutoff-cancellation} they may be assigned to overlapping
unselected support indicators.

\paragraph{Geometric derivatives.}
Equation~\eqref{can:geometry} gives $|r'|\lesssim1$ and
$|n'|\lesssim\theta/r$. Since $n'\perp n$,
\[
 \frac{|d'_n|}{d}\lesssim r^{-1},\qquad
 \frac{|D'_n|}{D}\lesssim\frac{\theta}{r\phi}\lesssim r^{-1}.
\]
Together with the logarithmic rate of $r$, these give
\begin{equation}\label{can:spatial-cutoff-bound}
 C\frac{f\phi}{r^2\theta}.
\end{equation}
These errors can also be assigned to unselected neighbors by
Lemma~\ref{can:cutoff-cancellation}.

\paragraph{Receiver derivatives.}
Only $D$ in the cutoff depends directly on $a$. Its source-time rate is
\[
 \frac{|D'_a|}{D}
   =\frac{|n\cdot a_t|}{D}\frac dD
   \lesssim\frac{\theta^2|K|}{w\phi^3}.
\]
Multiplication by $fDk_0/(rd)$ gives
\begin{equation}\label{can:receiver-cutoff-bound}
 C\frac{f\theta|K|}{wr\phi^2}.
\end{equation}
This is estimated on selected supports directly. Every bound
\eqref{can:source-cutoff-bound}--\eqref{can:receiver-cutoff-bound} gains
a factor $\phi$ after projection by $P_X$.

\paragraph{Time weights and a varying projection.}
If $\psi(t)$ is a Lipschitz weight, its derivative contributes
\begin{equation}\label{can:time-weight-bound}
 \psi_t\chi\,\frac{fk_0}{r},
 \qquad\left|\frac{fk_0}{r}\right|\lesssim\frac{f\theta}{r\phi}.
\end{equation}
For $|V_X|\asymp w$, put $\Pi=wP_X/|V_X|$. Its norm is bounded and
$\|\Pi_t\|\lesssim |K|/w$, by differentiation of the unit direction
$V_X/|V_X|$ and of $|V_X|^{-1}$. Consequently
\begin{equation}\label{can:projector-derivative}
 \left|\frac{f\Pi_tk_0}{r}\right|
       \lesssim\frac{f\theta|K|}{wr\phi}.
\end{equation}
This has the size of the projected receiver term. The estimate uses
$|k_0|\lesssim\theta/\phi$, so it retains the longitudinal part of $k_0$.

\subsection{Integration by parts and the initial cone}

Write $\mathcal K_t=fD\mathcal K$ for the interior force density in
$dt\,dy\,dv$ measure; the subscript identifies the measure and does not
denote a derivative.

\begin{proposition}[Weighted signed formula]\label{can:weighted-formula}
Let $J=[t_1,t_2]$ lie in a compact classical time interval. Let
$0\leq\psi\leq1$ be Lipschitz on $J$, vanish at its endpoints, and be
extended by zero outside $J$. Let $\chi$ be a selected sum as in
Lemma~\ref{can:cutoff-cancellation}. One may integrate
\eqref{can:label-identity} labelwise against $\psi\chi$ and a bounded
$C^1$ matrix $\Pi(t)$, including $\Pi=\Id$. More precisely,
\begin{align}\label{can:weighted-terms}
 \int_J\!\int\psi\chi\Pi\mathcal K_t\,dy\,dv\,dt
 ={}&\mathcal B_0+\int_J\!\int
 \biggl[\psi\chi\Pi(\mathcal R+\mathcal A)
       +\psi_t\chi\Pi\frac{fk_0}{r}\notag\\
 &\hspace{25mm}+\psi\chi\Pi_t\frac{fk_0}{r}
       +\psi\Pi\frac{fDk_0}{rd}\chi'_s\biggr]\,dy\,dv\,dt,
\end{align}
where the spatial integral has $r=|X(t)-y|<t$, all source quantities are
evaluated at $s=t-r$, and $\mathcal B_0$ is the initial-cone boundary term.
If $\|\Pi\|\lesssim1$, that
boundary term is bounded in norm by
\begin{equation}\label{can:initial-boundary-bound}
 C\int_J|\psi(t_0)|\,dt_0.
\end{equation}
The constant in \eqref{can:initial-boundary-bound} depends only on the
datum and fixed time horizon. For $\Pi=\Id$, and for
$\Pi=wP_X/|V_X|$ when $|V_X|\asymp w$, the derivative kernels satisfy
the bounds above, uniformly in the fixed selected index set. For a general
$\Pi$, its derivative term retains the explicit bound
$Cf\theta\|\Pi_t\|/(r\phi)$.
\end{proposition}

\begin{proof}
The null set of labels with a same-time collision with the receiver can
be discarded by Lemma~\ref{ret:receiver-time}. For any remaining
label, continuity on the compact time interval gives
$\min_t|Y(t)-X(t)|>0$. At a retarded point,
\[
 |Y(t)-X(t)|\leq |Y(t)-Y(s)|+r\leq2r,
\]
so $r$ is bounded below on that label's branch. The branch is $C^1$ and
increasing. Its entry at source time $s=0$, when present, has receiver
time $t_0$ determined by $t_0=|X(t_0)-Y(0)|$. All other integration
endpoints occur where $\psi=0$. Products with Lipschitz $\psi$ are
absolutely continuous along the branch. Integrating the total derivative
in \eqref{can:label-identity} is therefore legitimate for each such label.
The factors $\psi_s=\psi_t d/D$ and $\Pi_s=\Pi_t d/D$, followed by
\eqref{can:measure}, give exactly \eqref{can:weighted-terms}.

The sign at $s=0$ is positive. Its integrand relative to initial labels is
\[
 f_0(z)\psi(t_0)\chi\Pi(t_0)\frac{k_0}{t_0d(0)}.
\]
For fixed initial velocity, parameterize the initial position by
$y_{\rm in}=X(t_0)-t_0n$. This map is one-to-one on the initial cone
and has absolute Jacobian $t_0^2D$. The initial momentum support bounds
$d(0)$ below, and
\[
 Dk_0=Du-ne
\]
is bounded in norm. Since $\chi\leq1$, boundedness and compact velocity
support of $f_0$, and $t_0\leq T_0$, give
\eqref{can:initial-boundary-bound}. The same proof allows disjoint time
weights to be summed, using their sum at the initial boundary.

Finally these operations also hold after integrating the labels and
summing the dyadic products. On a fixed compact classical interval,
$d,D$ have positive lower bounds on the relevant supported trajectories,
and source and receiver accelerations are bounded. After conversion to
$dt\,dy\,dv$, all differentiated kernels have an integrable majorant
$C_{\rm comp} f(r^{-2}+r^{-1})$, using bounded partition overlap.
Indeed $dy=r^2dr\,dn$ and the momentum support is bounded. The constant
here may depend on compact classical bounds and on the fixed weight;
it is used only to justify the identity, not in the uniform estimates.
For a completely truncated argument, insert a radial cutoff vanishing for
$r<\gamma$ and equal to one for $r>2\gamma$. Its new derivative term is
bounded by $C_{\rm comp}f/(\gamma r)$ on $r\asymp\gamma$, whose integral
is $O_{\rm comp}(\gamma)$. All remaining terms converge by domination.
Local finiteness for $r>0$ and the same majorants justify the infinite sum
of distance bins and its differentiated version. Thus no cone-tip
boundary remains.
\end{proof}

\subsection{Coefficients used in the later selections}

We record a coefficient bound used in both selections, so that their
receiver-derivative costs can be read without repeating the geometry.
At a fixed receiver time, the coefficient on $|K|$ in the receiver term of
\eqref{can:bulk-bounds} or in \eqref{can:receiver-cutoff-bound}, integrated
over one small bin, obeys
\begin{equation}\label{can:receiver-coefficient}
 \frac Cw\min\left(h^2(p\theta)^3,
                       \frac1{p\theta h\phi^2}\right).
\end{equation}
For the first bound use the density estimate
$\int_{\rm bin}f\,dv\lesssim p^3\theta^2$, the $n$-cap area
$O(\phi^2)$, and $dy=r^2ds\,dn$. For the second, divide the kernel
$f\theta/(wr\phi^2)$ by the particle cone-flux density $qdf$ and use
the cone budget. Projection adds $\phi$ to the first bound; the
$\Pi_t$ term has the same projected size.

The direction estimate in Section~\ref{ang:section} will use the first,
projected coefficient with its own selection of indices; that argument
also accounts for the varying projector. The final signed-increment
estimate in Section~\ref{sel:section} will use the two-sided minimum
in \eqref{can:receiver-coefficient}. Both applications keep their
selected index sets fixed while differentiating the cutoffs.
